\enabstracttitle

With the increasing penetration of renewable energy and the deep peak-regulation operation of coal-fired power units, steam turbine systems are increasingly operated under wide-load and frequently varying conditions. Their pressure, temperature and flow parameters are affected not only by equipment health conditions, but also by unit load, steam boundaries and operating adjustments. Meanwhile, the steam turbine system exhibits strong coupling among multiple devices and parameters, whereas high-quality field fault samples are scarce. Consequently, conventional fixed-threshold methods and purely data-driven methods relying on labeled fault samples have difficulty simultaneously achieving adaptability to varying operating conditions, accurate fault identification and diagnostic interpretability. This thesis takes the thermal system of a coal-fired steam turbine as the research object and investigates equipment health condition monitoring, fault knowledge graph construction and KG-GCN-based fault diagnosis. The main research contents are as follows.

To address the continuous migration of normal steam turbine parameters with operating conditions, which makes a fixed design reference unsuitable for long-term condition assessment, a device-level health condition monitoring method based on a gated recurrent unit (GRU) is developed. According to the actual thermal boundary of each device, operating variables and upstream/external boundary variables are used as model inputs, while device-body and interface state parameters are used as outputs. The temporal model learns the healthy response under given operating conditions. The very-high-pressure turbine (VHP) is selected as a representative object. Normal DCS data from May 2024 are used for training, and data from June 2024 are used for independent testing, with ANN, TCN and iTransformer-small employed as comparison models. The results show that the GRU achieves a macro-average $R^2$ of 0.9511 over the final 24 state parameters and provides stable dynamic health references for the VHP body, first extraction-steam branch and first cold-reheat interface. These references provide the basis for representing equipment-state changes using residuals between measured values and healthy predictions.

To address the complex equipment relationships and the dispersion of fault knowledge across DCS point lists, equipment structure information and fault-related literature, a steam turbine fault knowledge graph is constructed from operation and technical materials. According to the actual thermal processes of the main steam path, reheating system, regenerative extraction-steam system and condenser cold end, equipment, parameter, anomaly, fault and maintenance-action entities together with their relation types are defined. A large language model is employed to generate candidate triples, which are subsequently checked and standardized using entity-type rules, relation-direction constraints, equipment compatibility, source evidence and embedding-based semantic alignment. The current audited knowledge base contains 127 standardized entities and 123 candidate relations, of which 121 relations are included in the main knowledge graph and two conflicting candidates are retained for further verification. In this way, distributed equipment-structure information and fault mechanisms are transformed into structured domain knowledge suitable for graph convolution computation.

To address fault-class confusion caused by shared local monitoring parameters and the limited availability of labeled fault samples, a fault diagnosis method based on a knowledge graph and graph convolutional network (KG-GCN) is proposed. Multi-parameter residuals generated by the health condition model are organized into temporal windows and mapped to parameter nodes of the knowledge graph, while parameter relations projected from the graph are used as graph-structure inputs. The GCN progressively aggregates fault features along equipment and parameter relations. During training, random masked-reconstruction pretraining is first performed on training residual windows without using class labels. The graph convolutional encoder is then frozen, and the classification layer is trained using labeled fault windows. Using the VHP as a representative case, four semi-physical diagnostic scenarios are constructed on real healthy residual backgrounds with independently defined fault mechanisms: normal condition, VHP flow-path erosion, first extraction-steam branch leakage and VHP seal leakage. In the fixed-seed experiment, KG-GCN achieves an Accuracy of 98.57\% and an F1 score of 98.58\%, higher than those of CNN and AE. Across three random seeds, its mean Macro-F1 reaches 98.46\%$\pm$0.11\%. When only 20\% of labeled fault windows are used in the classification stage, masked pretraining increases Macro-F1 from 94.48\% to 96.67\%, indicating improved diagnostic performance and stability under the current limited-label setting.

In summary, this thesis establishes a steam turbine condition monitoring and fault diagnosis framework consisting of device-level health modeling, dynamic state residual construction, steam turbine fault knowledge graph construction, graph-convolutional feature aggregation, masked pretraining and fault classification. The fault-classification experiments are based on real healthy residual backgrounds and semi-physical fault scenarios constructed from independent mechanisms. The results show that the proposed method provides good multi-class fault-identification performance in the current representative semi-physical case. Further validation using field fault records and simulation data calibrated against measurements is still required to assess engineering generalization to other steam turbine devices and fault conditions.

\enkeywords{Coal-fired power unit; Steam turbine system; Condition monitoring; Knowledge graph; Graph convolutional network; Fault diagnosis}
